% Sans protection, quatre dipeptides se forment (exercice 10, chapitre 12)
\begin{tikzpicture}[x=1cm,y=1cm,
  aa/.style={draw=otNavy,fill=white,line width=0.6pt,rounded corners=1.2mm,minimum width=15mm,minimum height=8mm,font=\otMedium\fontsize{8.5}{10}\selectfont,text=otInk},
  dp/.style={draw=otMuted,fill=otRule!25,line width=0.5pt,rounded corners=1mm,minimum width=22mm,minimum height=7mm,font=\otLight\fontsize{8}{9}\selectfont,text=otInk}]
\node[aa] (G) at (0,1.6) {Gly};
\node[aa] (A) at (0,0.3) {Ala};
\node[pcLegende,anchor=south] at (0,2.15) {chaque molécule porte \ce{-NH2} et \ce{-COOH}};
\draw[pcFleche] (1.6,0.95) -- (2.9,0.95);
\node[dp] at (4.6,1.95) {Gly--Gly};
\node[dp] at (4.6,1.05) {Gly--Ala};
\node[dp] at (4.6,0.15) {Ala--Gly};
\node[dp] at (4.6,-0.75) {Ala--Ala};
\node[pcLegende,anchor=west,align=left,text width=46mm] at (6.3,0.6) {en mélange direct, tous les couplages\\sont possibles : quatre dipeptides,\\dont un seul est celui recherché};
\end{tikzpicture}
\par\smallskip\legende{Le mélange brut des deux acides aminés conduit aux quatre enchaînements possibles.}
